\BourbakiAppendixExercises{习题}

\begin{enumerate}
\def\labelenumi{\arabic{enumi})}
\tightlist
\item
  设 \(\mathrm{A}\) 是一个 \(k\)-代数。
\end{enumerate}

\begin{enumerate}
\def\labelenumi{\alph{enumi})}
\item
  证明：包含 AA 的正则（左或右）理想只有 A 本身。
\item
  设 \(\mathfrak{a}\) 是 \(A\) 的一个双边理想。它既作为左理想又作为右理想都是正则的，当且仅当 \(k\)-代数 \(A / \mathfrak{a}\) 有单位元。
\end{enumerate}

\begin{enumerate}
\def\labelenumi{\arabic{enumi})}
\setcounter{enumi}{1}
\tightlist
\item
  设 \(\mathbf{A}\) 是一个 \(k\)-代数，\(\mathfrak{a}\) 与 \(\mathfrak{b}\) 是 \(\mathbf{A}\) 的正则（左）理想。
\end{enumerate}

\begin{enumerate}
\def\labelenumi{\alph{enumi})}
\item
  理想 \(a \cap b\) 是正则的，当且仅当存在模 a 的右单位元 u 与模 b 的右单位元 v，使得 \(u - v \in a + b\)。
\item
  由此推出：若 a 是双边理想且是右正则的，或者若 a 是极大的，则 \(a \cap b\) 是正则的。c) 设 K 是一个交换域，L 是两个变元 X 与 Y 上的自由结合 K-代数（III, §2, No.~7, 第 448 页），而 A 是 L 中由次数 \(\geqslant 1\) 的元素组成的 K-子代数。证明 A 的理想 \(A(1 - X)\) 与 \(A(1 - Y)\) 都是正则的，但它们的交约化为 0。
\end{enumerate}

\begin{enumerate}
\def\labelenumi{\arabic{enumi})}
\setcounter{enumi}{2}
\tightlist
\item
  设 \(\mathbf{A}\) 是一个 \(k\)-代数，\(\mathfrak{a}\) 是 \(\mathbf{A}\) 的一个正则左理想，\(u\) 与 \(v\) 是模 \(\mathfrak{a}\) 的右单位元。\(a)\) 给出一个 \(u - v\) 不属于 \(\mathfrak{a}\) 的例子（参见 III, §9, 第 644 页，习题 1）。由此推出一个使映射 \(\widetilde{\mathfrak{a}} \mapsto \widetilde{\mathfrak{a}} \cap \mathbf{A}\)（VIII, 第 437 页，命题 3）不是单射的例子。
\end{enumerate}

\begin{enumerate}
\def\labelenumi{\alph{enumi})}
\setcounter{enumi}{1}
\tightlist
\item
  假设理想 \(\mathfrak{a}\) 是双边理想，且 \(k\)-代数 \(B = A / \mathfrak{a}\) 不包含任何非零双边理想 \(\mathfrak{b}\) 使得 \(B\mathfrak{b} = 0\)。证明 \(u - v\) 属于 \(\mathfrak{a}\)。特别地，当 \(A\) 交换时就是这种情形。
\end{enumerate}

\begin{enumerate}
\def\labelenumi{\arabic{enumi})}
\setcounter{enumi}{3}
\tightlist
\item
  设 A 是一个 \(k\)-代数，\(\mathfrak{a}\) 是 A 的一个左理想；我们把 \(\mathfrak{a}\) 视为一个 \(k\)-代数。证明我们有 \(\Re(A) \cap \mathfrak{a} \subset \Re(\mathfrak{a})\)（应用 VIII, 第 441 页，命题 7），并且当 \(\mathfrak{a}\) 是双边理想时有等式（注意 A 上的每个单伪模在 \(\mathfrak{a}\) 上是单的或被 \(\mathfrak{a}\) 零化）。给出一个左理想 \(\mathfrak{a}\) 的例子，使得 \(\Re(A) \cap \mathfrak{a} \neq \Re(\mathfrak{a})\)（取 A 为一个矩阵代数）。
\end{enumerate}

¶ 5) 假设 A 上的每个伪模 M 都是 AM 与 M 中由被 A 零化的元素组成的子伪模的直和。证明 A 有单位元。（取 \(M = \widetilde{A}_{s}\)，证明 A 容许一个右单位元 u。取 \(M = A_{s}\)，证明对 A 的每个非零元素 x 我们有 \(Ax \neq 0\)。由此断定 u 也是一个左单位元。）

\begin{enumerate}
\def\labelenumi{\arabic{enumi})}
\setcounter{enumi}{5}
\tightlist
\item
  我们称一个 k-代数是本原的，如果它容许一个既单又忠实的伪模（即其零化子约化为 0 的伪模）。
\end{enumerate}

\begin{enumerate}
\def\labelenumi{\alph{enumi})}
\item
  证明：若代数 \(\tilde{\mathbf{A}}\) 是本原的，则 \(\mathbf{A}\) 也是本原的。
\item
  若代数 A 是本原的且有单位元，则代数 \(\widetilde{A}\) 不是本原的。
\item
  假设 k 是一个体，且代数 A 是本原的但没有单位元。证明 \(\widetilde{A}\) 是本原的（证明：若存在 A 上的一个忠实伪模，它作为 \(\widetilde{A}\)-模不忠实，则 A 容许一个右单位元 u，并由此推出矛盾）。
\end{enumerate}

\begin{enumerate}
\def\labelenumi{\arabic{enumi})}
\setcounter{enumi}{6}
\item
  我们称 \(k\)-代数 A 是阿廷的，如果 A 的左理想的每个递减序列都是平稳的。环 \(\widetilde{\mathrm{A}}\) 是阿廷的，当且仅当代数 A 与环 \(k\) 都是阿廷的。证明：若 \(k\) 是一个体，则无根的阿廷代数 A 有单位元（注意代数 \(\widetilde{\mathrm{A}}\) 是半单的）。
\item
  假设 \(k\) 是一个体，且 \(A\) 是 \(k\) 上的一个本原代数。证明：若 \(A\) 容许一个在 \(k\) 上有限维的非零左理想，则代数 \(A\) 在 \(k\) 上是有限维的，从而是一个单环（参见习题 7）。
\item
  假设 k 是一个体，且 A 是一个幂零 k-代数，即存在一个整数 n 使得 \(A^{n}=0\)。证明 A 在 k 上是有限维的（对使 \(A^{n}=0\) 的最小整数 n 用归纳法）。
\item
  假设 \(k\) 是一个体，\(A\) 是一个无限维 \(k\)-代数，并且 \(A\) 的每个真左理想都在 \(k\) 上有限维。
\end{enumerate}

\begin{enumerate}
\def\labelenumi{\alph{enumi})}
\item
  证明我们有 \(Aa = 0\) 对每个左理想 \(a \neq A\) 成立（考虑同态 \(\gamma : A \to \operatorname{End}_{k}(a)\)）。
\item
  设 \(\mathfrak{r}\) 是 A 的极大左理想之和。证明 \(\mathfrak{r}\) 异于 A（利用习题 9），\(\mathfrak{r}\) 是一个双边理想，且 \(A / \mathfrak{r}\) 是一个体（参见习题 7）。c) 设 \(u\) 是模 \(\mathfrak{r}\) 的一个（左或右）单位元。证明我们有 \(\mathfrak{ru} = \mathfrak{r}\)（注意我们有 \(Au = A\)）。
\item
  证明 rA 为零（考虑 r 在 A 中的右零化子，并按 a) 的方式推理），然后证明 r 为零。由此得出 A 是一个体。
\end{enumerate}

¶ 11) 假设 k 是一个体，A 是一个 k-代数，并且 A 的子代数的每个递减序列都是平稳的。

\begin{enumerate}
\def\labelenumi{\alph{enumi})}
\item
  证明代数 \(\mathbf{A}\) 是阿廷的，且其根在 \(k\) 上有限维（应用习题 9）。
\item
  设 S 是 A/ \(\Re(A)\) 的一个单分量（习题 7）；它同构于一个体 D 上的矩阵代数，其中心为 Z。证明 D 在 Z 上有限维（若 D 含有一个在 Z 上超越的元素，构造 D 中一个无穷严格递减的子体序列；若 D 在 Z 上是代数的且无限维，构造 D 中一个无穷严格递增的子体序列，并考虑它们在 D 中的交换子）。
\item
  证明：若 Z 在 k 上是无限维的，则我们有 S = D（考虑 B 的幂零子代数）。
\item
  证明 Z 是 k 的一个代数扩张。设 p 是 k 的特征指数，L 是 Z 所含的 k 的最大可分扩张。证明 \(L \cap Z^{p}\) 在 \(L^{p}\) 上有有限次数（考虑 \(L \cap Z^{p}\) 在 \(L^{p}\) 上的一个 p-基）。
\item
  L 的子扩张的每个递减序列都是平稳的。给出一个例子说明这个条件并不蕴含 L 在 k 上有有限次数（参见 V, §12, 第 167 页，习题 3）。
\item
  反之，设 B 是一个满足 a) 的条件的 k-代数，使得 \(\mathrm{B}/\Re(\mathrm{B})\) 的每个单分量都满足 b)、c)、d) 与 e) 的条件。证明 B 的子代数的每个递减序列都是平稳的。
\end{enumerate}

（归结为 B 是一个体且中心为 Z 的情形。设 \((\mathrm{D}_{n})\) 是 B 的子体的一个递减序列；我们可以假设由 \(D_{n}\) 与 Z 生成的体 D 与 n 无关。注意我们可以写 \(D = E \otimes_{L} Z\)，其中 L 是 Z 的一个在 k 上有有限次数的子体，而 E 是一个在 L 上有有限次数的体，其中心包含 L。证明 \(D_{n}\) 这时就是由 E 与 \(D_{n} \cap Z\) 生成的体。于是我们归结为 B 是一个交换域且是 k 的 p-根式扩张的情形。若 \((\mathrm{F}_{n})\) 是 B 的子扩张的一个递减序列，考虑子体 \(k(\mathrm{F}_{n}^{p^{\infty}})\)，其中 \(F_{n}^{p^{\infty}} = \cap_{f} F_{n}^{p^{f}}\)。）

¶ 12) a) 假设 k 是一个体，k-代数 A 是阿廷的，并且 A 的子代数的每个递增序列都是平稳的。证明 A 在 k 上有有限次数。（按习题 11 的方式推理。注意体 \(k(\mathrm{X})\) 含有一个无穷严格递增的子代数序列，这利用了 VII, §1, No.~5, 第 6 页，命题 6。）b) 证明在多项式代数 \(k[X]\) 中，子代数的每个递增序列都是平稳的。（设 \(f \in k[X]\) 是一个非常数多项式，A 是由 1 与 f 生成的 \(k[X]\) 的 k-子代数。注意 A 是一个主理想整环，且 \(k[X]\) 是一个有限生成 A-模。）

\begin{enumerate}
\def\labelenumi{\arabic{enumi})}
\setcounter{enumi}{12}
\item
  假设 \(k\) 是一个体，且 \(A\) 容许一个由幂零元素组成的在 \(k\) 上的有限基。证明 \(A\) 是幂零的（习题 9；归结为代数 \(A\) 是半单的情形（参见习题 7），然后通过扩张基域归结为形如 \(\mathbf{M}_n(\mathrm{K})\) 的情形）。
\item
  设 A 与 B 是 k-代数；(A,B)-伪双模是一个集合 E，它赋有 A 上的左伪模结构与 B 上的右伪模结构，且使得底 k-模结构一致。这对应于赋予 E 一个 \(\widetilde{A}\) 与 \(\widetilde{B}\) 这两个代数上的双模结构。E 的一个乘子是一个偶对 (S,D)，其中 S:B→E 是右 B-伪模的同态，D:A→E 是左 A-伪模的同态，满足 \(a\mathrm{S}(b)=\mathrm{D}(a)b\)（对 \(a\in A\)、\(b\in B\)）。用 \(\mathcal{M}(\mathrm{E})\) 表示 E 的乘子所成的集合；它容许一个自然的 k-模结构。
\end{enumerate}

\(a\) ) 对每个 \(x\in \mathrm{E}\)，用 \(\mathrm{D}_x\)（相应地，\(\mathrm{S}_x\)）表示映射 \(a\mapsto ax\)（相应地，\(b\mapsto xb\)）。证明偶对 \(m_{x} = (\mathrm{D}_{x},\mathrm{S}_{x})\) 是 E 的一个乘子。

\begin{enumerate}
\def\labelenumi{\alph{enumi})}
\setcounter{enumi}{1}
\tightlist
\item
  对每个 \(m = (\mathrm{S}, \mathrm{D}) \in \mathcal{M}(\mathrm{E})\)、\(a \in A\) 与 \(b \in B\)，置 \(am = m_{\mathrm{D}(a)}\) 与 \(mb = m_{\mathrm{S}(b)}\)。证明这在 k-模 \(\mathcal{M}(\mathrm{E})\) 上定义了一个 (A, B)-伪双模结构，且映射 \(x \mapsto m_x\) 是 (A, B)-伪双模的一个同态。
\end{enumerate}

\begin{enumerate}
\def\labelenumi{\arabic{enumi})}
\setcounter{enumi}{14}
\tightlist
\item
  设 \(\mathrm{A}\) 是一个 \(k\)-代数。\(\mathrm{A}\) 的乘子是 (A,A)-伪双模 \(\mathrm{A}\) 的一个乘子。乘子所成的 \(k\)-代数（\(\mathrm{A}\) 的）用 \(\mathcal{M}(\mathrm{A})\) 表示。
\end{enumerate}

\begin{enumerate}
\def\labelenumi{\alph{enumi})}
\item
  证明映射 \(\mu: a \mapsto m_{a}\) 是一个 k-代数同态，且我们有 \(am = m_{a}m\) 与 \(ma = mm_{a}\)（对 \(a \in A, m \in \mathcal{M}(A)\)）。因此同态 \(\mu\) 的像是环 \(\mathcal{M}(A)\) 的一个双边理想。
\item
  k-代数 A 的一个双边理想 b 称为忠实的，如果它的左零化子与右零化子都为零。证明：若 b 在 A 中是忠实的，则同态 \(\mu\) 是单射，且理想 \(\mu(\mathrm{A})\) 在 \(\mathcal{M}(\mathrm{A})\) 中是忠实的。
\item
  设 B 是一个包含 A 作为双边理想的有单位元的 k-代数。对 \(b \in B\)，设 \(\sigma_{b}\) 与 \(\delta_{b}\) 是从 A 到 A 的映射，使得对 \(a \in A\) 有 \(\sigma_{b}(a) = ba\) 与 \(\delta_{b}(a) = ab\)，并置 \(\pi(b) = (\sigma_{b}, \delta_{b})\)。证明 \(\pi(b) \in \mathcal{M}(A)\) 对每个 \(b \in B\) 成立，且如此得到的映射 \(\pi : B \to \mathcal{M}(A)\) 是一个环同态，它在 A 上的限制是映射 \(a \mapsto m_{a}\)。证明：若理想 A 在 B 中是忠实的，则 \(\pi\) 是单射。
\end{enumerate}

